\subsection{What changes the outcome: information or learning design?}
\label{sec:design}

The two channels make opposite predictions about interventions. If low trading
were sustained by punishment, degrading the information traders use to
monitor each other should restore competition, and the content of their memory
should matter. If it is pruning, information is irrelevant and changes to the
learning procedure should matter. Table~\ref{tab:design} reports treatments
in the standard Kyle market, each against the residual-memory baseline
($\Delta=\ResidDeltaInt$).

\paragraph{Information does not matter.} Blurring the order flow that traders
observe with noise of twice the noise-trading volume leaves the outcome
unchanged ($\Delta=\OpaqueDeltaInt$). Replacing the informative memory with a
\emph{random} state that carries no information at all, but has the same
number of states, reproduces it ($\Delta=\RandomMemDeltaInt$). The difference
between forward-looking learners with and without memory
(\ResidDeltaInt{} versus \NoneDeltaInt{}) is therefore a matter of table size:
more states mean fewer updates per entry and more room for the selection effect
of Section~\ref{sec:mechanism}, not more information to coordinate on.

\paragraph{Learning design matters.} Two changes to the learning rule move the
outcome substantially. Counterfactual updating, which removes the selection
effect, eliminates the supracompetitive outcome: forward-looking learners end
at $\Delta=\CfResidDeltaInt$ with residual memory and $\CfNoneDeltaInt$ without,
that is, at (marginally beyond) the competitive level, with price
informativeness at its Nash value ($\Delta$ informativeness
$\CfResidDeltaInfo$). The market, the agents' information and the exploration
schedule are unchanged; only the update rule differs. Per-entry decaying step sizes
change even the sign: with the schedule of Section~\ref{sec:method}, learners
\emph{over}-trade relative to Nash ($\Delta=\VisitsResidual$ with residual
memory, $\VisitsNone$ without), because step sizes shrink while the market
maker's price impact is still depressed by heavy exploration, freezing
estimates around best responses to a market that is too cheap to trade in.
These learners are not converged either (on-path orders still move by
$\VisitsShift$ grid steps), which underlines how strongly the learning
procedure, rather than the game, shapes the outcome.

\paragraph{Market structure.} Adding two passive informed traders who play the
Nash strategy changes the benchmarks (the cartel now faces outside
competition) and yields $\Delta=\PassiveDeltaInt$. Without punishment in any
treatment (rival reactions in Table~\ref{tab:design} are indistinguishable from
zero and deviations pay), we read this as a change in what the biased learners
converge to rather than as a disruption of collusion.

\begin{table}[t]
  \centering
  \small
  \resizebox{\linewidth}{!}{\begin{tabular}{lcccc}
\toprule
Treatment ($\xi=0$, $I=2$) & $\Delta$ intensity & $\Delta$ profit & Rival $\Delta\beta_1$ & Deviator gain \\
\midrule
Random memory, 35 states & $0.84 \pm 0.01$ & $0.45 \pm 0.03$ & $0.000 \pm 0.000$ & $0.052 \pm 0.004$ \\
Random memory, myopic & $1.16 \pm 0.01$ & $0.11 \pm 0.02$ & $0.000 \pm 0.000$ & $0.107 \pm 0.006$ \\
Counterfactual updates, residual memory & $-0.10 \pm 0.01$ & $-0.44 \pm 0.03$ & $-0.001 \pm 0.002$ & $0.020 \pm 0.003$ \\
Counterfactual updates, no memory & $-0.09 \pm 0.03$ & $-0.34 \pm 0.08$ & $0.000 \pm 0.000$ & $0.015 \pm 0.002$ \\
Reduced transparency (disclosure noise $2\sigma_u$) & $0.90 \pm 0.01$ & $0.46 \pm 0.03$ & $-0.002 \pm 0.003$ & $0.058 \pm 0.004$ \\
\bottomrule
\end{tabular}}
  \caption{Treatments in the standard Kyle market ($\xi=0$, $I=2$, residual
  memory unless noted, $6\times10^6$ periods, 200 markets). Baseline:
  $\Delta$ intensity $\ResidDeltaInt$. 95\% CIs.}
  \label{tab:design}
\end{table}
